% A two-column conference paper. The preamble (up to \begin{document}) gives
% it its look; the writing starts at the abstract.
\documentclass[10pt,letterpaper]{article}

% Fonts and typography
\usepackage[T1]{fontenc}
\usepackage{mathptmx}
\usepackage{microtype}

% Page: two columns
\usepackage[letterpaper,hmargin=0.75in,top=0.85in,bottom=1in,columnsep=0.3in]{geometry}

% Mathematics, figures and tables
\usepackage{amsmath,amssymb}
\usepackage{graphicx}
\graphicspath{{figures/}}
\usepackage{booktabs}
\setlength{\belowcaptionskip}{6pt} % tables carry their caption on top

% Citations, numeric (natbib with BibTeX)
\usepackage[numbers,sort&compress]{natbib}

% Links and citations in the paper's accent colour
\usepackage{xcolor}
\definecolor{accent}{RGB}{140,82,18}
\usepackage[colorlinks=true,linkcolor=accent,citecolor=accent,urlcolor=accent]{hyperref}

% Compact section headings for two columns
\makeatletter
\renewcommand\section{\@startsection{section}{1}{\z@}%
  {-2.2ex \@plus -0.5ex \@minus -0.2ex}{1ex \@plus 0.2ex}%
  {\normalfont\large\bfseries}}
\renewcommand\subsection{\@startsection{subsection}{2}{\z@}%
  {-1.8ex \@plus -0.5ex \@minus -0.2ex}{0.8ex \@plus 0.2ex}%
  {\normalfont\normalsize\bfseries}}
\makeatother

\newcommand{\papertitle}{{{title}}}
\newcommand{\paperauthor}{{{author}}}
\newcommand{\paperaffiliation}{Your institution}
\newcommand{\paperdate}{{{date}}}
\newcommand{\paperkeywords}{typesetting, literate programming, writing tools}

\hypersetup{pdftitle={\papertitle},pdfauthor={\paperauthor}}

\begin{document}

\twocolumn[
  \begin{center}
    {\LARGE\bfseries \papertitle\par}
    \vspace{0.9em}
    {\large \paperauthor\par}
    \vspace{0.2em}
    {\itshape \paperaffiliation\par}
    \vspace{0.2em}
    {\small \paperdate\par}
  \end{center}
  \vspace{1.5em}
]

\noindent\textbf{\textit{Abstract}---}One paragraph on the problem, the
approach, and the main result, with the number that matters. Write it last.

\smallskip
\noindent\textbf{\textit{Keywords}---}\paperkeywords

\section{Introduction}
\label{sec:introduction}

State the problem and why it matters in the first paragraph. Programs, like
papers, should be written for people first~\citep{knuth1984}; so should the
opening of a paper.

End the introduction with the contributions:
\begin{itemize}
  \item the first contribution, in one line;
  \item the second, with the section that shows it
    (Section~\ref{sec:approach}).
\end{itemize}

\section{Approach}
\label{sec:approach}

Describe the method precisely enough to reproduce it. Equations are numbered:
\begin{equation}
  \hat{y} = \operatorname*{arg\,max}_{y}\, p_\theta(y \mid x)
  \label{eq:decode}
\end{equation}
and referred to by label, as in~\eqref{eq:decode}.

\section{Evaluation}
\label{sec:evaluation}

\begin{table}[htbp]
  \centering
  \caption{Results on the main benchmark.}
  \label{tab:main}
  \begin{tabular}{lrr}
    \toprule
    Method   & Score & Time (s) \\
    \midrule
    Method A & 71.2  & 3.4      \\
    Method B & 84.9  & 2.1      \\
    \bottomrule
  \end{tabular}
\end{table}

Table~\ref{tab:main} compares the methods. Put images in the figures
folder and place them in a \texttt{figure} environment with a caption.

\section{Related Work}
\label{sec:related}

Place the paper among the closest prior work, citing each by its key.

\section{Conclusion}
\label{sec:conclusion}

Summarize the result and its limits in a few sentences.

\bibliographystyle{abbrvnat}
\bibliography{refs}

\end{document}
